\itemhead{p.84 练习 · 第4题（教材p.84）}{XJ-B1-C03-PR-P084-Y1532}
\label{XJ-B1-C03-PR-P084-Y1532}
{\footnotesize 湘教版教材p.84 / PDF第91页 / p.84 练习 / 第4题}
\field{原题}{设 $g(x)$ 是定义于 $(-\infty,+\infty)$ 上的函数，且 $f(x)=g(x)-g(-x)$，讨论 $f(x)$ 的奇偶性；如果在 $[0,+\infty)$ 上 $f(x)=x^2-x$，求 $f(x)$ 在 $(-\infty,0]$ 上的解析式。}
\field{参考答案}{$f(-x)=g(-x)-g(x)=-f(x)$，为奇函数；$x\le0$ 时 $f(x)=-f(-x)=-x^2-x$，在0处也与原式一致。}
{\small \field{关联知识}{F08 奇偶性、对称与函数性质应用（教材帮p.113、114、115 / PDF122、123、124）；F03 函数求值、解析式与函数建模（教材帮p.80、87、88、89、90、131 / PDF89、96、97、98、99、140）}
\field{方法与易错点}{检查定义域，再代入 $-x$；求另一侧解析式时同时写出 $-x$ 所在区间。 易错：只看公式不看定义域；用一对点证明奇偶；负半轴代换遗漏最外层负号。}
\field{答案依据}{答案册相关题组：PDF第23页（册内页22）；按当前题库原题逐项重算，不按旧版题号机械对套。}}
